% Claim statement file for row claim_3_9 -- "SwigAcyclic".
% Auto-generated stub by scaffold/solve_chapter.py at row start. The
% manager agent (or a worker it dispatches) fills the body below with the
% claim's statement(s) from the lecture notes -- statement only, no proof
% (proofs live in the sibling `claim_3_9_proof_*.tex` file).
%
% This file is a sibling of `main.tex` inside `<subsection>/tex/`. The
% `subfiles` package lets it render both standalone and as part of
% `main.tex`. Note: `main.tex` does NOT \subfile this statement file -- the
% sibling `_proof_` file restates the statement above the proof, so
% including both in `main.tex` would be redundant. This file exists for
% standalone reading / grep convenience.

\documentclass[main]{subfiles}

\begin{document}
% ref: claim_3_9 (statement)
% title: SwigAcyclic
\def\rowref{claim\_3\_9}
\phantomsection\label{claim_3_9}

\begin{Rem}[SwigAcyclic]
    Let $G = (J, V, E, L)$ be a CADMG, i.e.\ a CDMG (def \ref{def-cdmg}) that is acyclic in the sense of def \ref{def-acylic} (equivalently, $G$ satisfies the attribute ``CADMG'' of def \ref{def-cdmg-types}, item~i.); throughout we use the notation of def \ref{not-cdmg}, the family-relationship notation of def \ref{def_3_5}, the topological-order definition of def \ref{def-topological-order}, and the SWIG (node-splitting hard-intervention) construction of def \ref{def:G_node-splitting_intervention}. Let
    \[ W \ins V \]
    be a subset of the output-node set of $G$, and let
    \[ G_{\swig(W)} \;=\; \lp J_{\swig(W)}, V_{\swig(W)}, E_{\swig(W)}, L_{\swig(W)} \rp \]
    be the SWIG of $G$ w.r.t.\ $W$ as constructed in def \ref{def:G_node-splitting_intervention}. Both quantifiers ``for every CADMG $G$'' and ``for every $W \ins V$'' are read \emph{universally}, with the side condition $W \ins V$ inherited verbatim from def \ref{def:G_node-splitting_intervention}; the case $W = \emptyset$ is admitted (then $W^o = W^i = \emptyset$ and $G_{\swig(W)} = G$ by items i.--iv.\ of def \ref{def:G_node-splitting_intervention}).

    \emph{Carrier of $G_{\swig(W)}$.} By items i.\ and ii.\ of def \ref{def:G_node-splitting_intervention},
    \[
        J_{\swig(W)} \cup V_{\swig(W)} \;=\; J \cup (V \sm W) \dcup W^o \dcup W^i,
    \]
    where $W^o := \lC w^o \st w \in W \rC$ and $W^i := \lC w^i \st w \in W \rC$ are the two tagged copies of $W$ introduced in def \ref{def:G_node-splitting_intervention}. The four constituents $J$, $V \sm W$, $W^o$, $W^i$ are pairwise type-level disjoint by def \ref{def:G_node-splitting_intervention}: $J \cap V = \emptyset$ (def \ref{def-cdmg}), $(V \sm W) \cap W = \emptyset$, and the tagged copies $W^o$, $W^i$ are formally distinct from each element of $J \cup V$ and from each other ($W^o \cap W^i = \emptyset$, $w^o \neq w^i$ for every $w \in W$).

    \emph{Finiteness and indexing of $J \cup V$.} By def \ref{def-cdmg} items~i.--ii., the set $J \cup V$ is finite; set
    \[ n \;:=\; |J \cup V|. \]

    \emph{The remark.} The following conjunction of sub-claims holds.

    \begin{enumerate}[label=(\alph*)]
        \item \emph{$G_{\swig(W)}$ is acyclic.}
              The CDMG $G_{\swig(W)}$ is acyclic in the sense of def \ref{def-acylic}, i.e.\ for every $x \in J_{\swig(W)} \cup V_{\swig(W)} = J \cup (V \sm W) \dcup W^o \dcup W^i$, there does not exist any non-trivial directed walk from $x$ to itself in $G_{\swig(W)}$.

        \item \emph{An explicit topological order on $G_{\swig(W)}$ obtained from any topological order on $G$.}
              For every topological order $<$ of $G$ in the sense of def \ref{def-topological-order}, presented in the equivalent indexing form (def \ref{def-topological-order}, ``Equivalent indexing form'') as a bijection
              \[
                  \{1, 2, \dots, n\} \;\longrightarrow\; J \cup V, \qquad i \;\longmapsto\; v_i,
              \]
              with $v_1 < v_2 < \cdots < v_n$ (so that ``parents precede children'' holds in this enumeration: $v_i \in \Pa^G(v_j) \Rightarrow i < j$), the binary relation $<_{\swig}$ on $J_{\swig(W)} \cup V_{\swig(W)}$ defined below is a topological order of $G_{\swig(W)}$ in the sense of def \ref{def-topological-order}.

              \emph{Index map.}
              Define the index map
              \[ \phi \,:\, J_{\swig(W)} \cup V_{\swig(W)} \;\longrightarrow\; \Q \]
              by case analysis on the disjoint-union carrier $J \cup (V \sm W) \dcup W^o \dcup W^i$:
              \begin{itemize}
                  \item \emph{Unsplit nodes retain their original index.}
                        For every $x \in J \cup (V \sm W)$, write $x = v_j$ with $j \in \{1, 2, \dots, n\}$ the unique index given by the enumeration of $J \cup V$ above, and put
                        \[ \phi(x) \;:=\; j. \]
                  \item \emph{Tagged copies of $W$ are placed symmetrically around the original index, with the output-side copy first.}
                        For every $w \in W$, write $w = v_j$ with $j \in \{1, 2, \dots, n\}$ the unique index given by the enumeration of $J \cup V$ above, and put
                        \[
                            \phi(w^o) \;:=\; j - \tfrac{1}{3}, \qquad \phi(w^i) \;:=\; j + \tfrac{1}{3}.
                        \]
              \end{itemize}

              \emph{Order from index map.}
              Define
              \[
                  x \;<_{\swig}\; y \quad :\Longleftrightarrow\quad \phi(x) \,<\, \phi(y) \qquad \text{for } x, y \in J_{\swig(W)} \cup V_{\swig(W)},
              \]
              where the right-hand inequality is the standard strict order on $\Q$.
    \end{enumerate}

    \emph{Where each clause of the LN's prose lands.}
    The LN's clauses ``$v_j^o$ the index $j - \tfrac{1}{3}$'' and ``$v_j^i$ the index $j + \tfrac{1}{3}$'' for $v_j \in W$ are the second bullet of the index map above. The LN's closing prose ``and then ordering all nodes according to their index value'' is the definition of $<_{\swig}$ from $\phi$ via the strict order on $\Q$. The first bullet of the index map -- ``unsplit nodes $v_j \in J \cup (V \sm W)$ retain index $j$'' -- is the rule the LN leaves implicit: the LN's literal text assigns indices only to the tagged copies $v_j^o$, $v_j^i$ of split nodes $v_j \in W$ and says nothing about the indices of $v_j \in J \cup (V \sm W)$, but the closing instruction ``order all nodes according to their index value'' is well-defined on the carrier $J_{\swig(W)} \cup V_{\swig(W)}$ only if the unsplit nodes also carry an index. The natural and only-consistent reading is that the original index $j$ is retained for unsplit nodes; the first bullet records this retention rule explicitly.

    \emph{Load-bearing role of the offset $\tfrac{1}{3}$.}
    The LN's specific numerical choice $\tfrac{1}{3}$ for the offset is not aesthetic: any value $\delta \in (0, \tfrac{1}{2})$ would equally well define a strict total order $<_{\swig}$ on $J_{\swig(W)} \cup V_{\swig(W)}$ realising the topological-order property, but $\delta = \tfrac{1}{2}$ already fails. Concretely, if $v_j \in W$ and $v_{j+1} \in W$ are two consecutively-indexed split nodes, then with $\delta = \tfrac{1}{2}$ the values $\phi(v_j^i) = j + \tfrac{1}{2}$ and $\phi(v_{j+1}^o) = (j+1) - \tfrac{1}{2} = j + \tfrac{1}{2}$ collide, so the relation ``order all nodes according to their index value'' fails to be a strict total order on the carrier. We adopt the LN's literal choice $\tfrac{1}{3}$ here; the role this choice plays in the proof of the topological-order property is deferred to the proof file.

    \emph{Dependence on the no-transfer-edge clause of def \ref{def:G_node-splitting_intervention}, item iii.}
    The index assignment $\phi(w^o) = j - \tfrac{1}{3} < j + \tfrac{1}{3} = \phi(w^i)$ places the output-side copy $w^o$ strictly before the input-side copy $w^i$ in $<_{\swig}$. Sub-claim~(b)'s topological-order assertion is consistent with this choice only because def \ref{def:G_node-splitting_intervention}, item~iii., introduces \emph{no} transfer edge in the $w^o \to w^i$ direction: the typing of item~iii.\ forces every directed edge of $G_{\swig(W)}$ to have its source tagged $\cdot^i$ (or in $J \cup (V \sm W)$ under the shorthand of def \ref{def:G_node-splitting_intervention}) and its target tagged $\cdot^o$ (or in $V \sm W$ under that shorthand); hence $(w^o, w^i) \notin E_{\swig(W)}$ by the structural typing alone, no matter what edges $E$ contains. The opposite direction $(w^i, w^o)$ is structurally permitted by item~iii.\ (it would lift from a directed self-loop $(w, w) \in E$ under the set-builder), but it is ruled out under the standing CADMG hypothesis on $G$ -- see the next paragraph. (Contrast with def \ref{def:G_node-splitting}, item~iii., which \emph{does} introduce a transfer edge $(w^0, w^1) \in E_{\spl(W)}$ for every $w \in W$ -- the analogous claim for the node-splitting construction, claim_3_6, has its topological-order assignment $\phi(w^0) = j - \tfrac{1}{3}$, $\phi(w^1) = j + \tfrac{1}{3}$ aligned precisely with the source/target slots of that transfer edge, so that $w^0 <_{\spl} w^1$ is forced by the topological-order property. For the SWIG, no transfer edge in either direction constrains the relative order of $w^o$ and $w^i$, and the choice $w^o <_{\swig} w^i$ is recorded by the index assignment itself.)

    \emph{Role of the CADMG (acyclicity) precondition.}
    The construction of def \ref{def:G_node-splitting_intervention} on its own does not produce a graph that satisfies sub-claims~(a) and~(b) for an arbitrary CDMG $G$; the standing CADMG hypothesis on $G$ (i.e.\ acyclicity of $G$ in the sense of def \ref{def-acylic}) is needed in two ways.
    \begin{itemize}
        \item For sub-claim~(a) (acyclicity of $G_{\swig(W)}$): a directed cycle $v_{i_1} \tuh v_{i_2} \tuh \cdots \tuh v_{i_k} \tuh v_{i_1}$ in $G$ whose nodes lie in $V \sm W$ (none of the cycle's nodes belonging to $W$) lifts under item~iii.\ to a directed cycle of the same length in $G_{\swig(W)}$, because the shorthand $v^i = v^o = v$ for $v \notin W$ makes each lifted edge $v_{i_a}^i \tuh v_{i_{a+1}}^o$ coincide with the original edge $v_{i_a} \tuh v_{i_{a+1}}$. Acyclicity of $G$ rules out all such cycles in $G$ outright.
        \item For sub-claim~(b) (the index assignment is a topological order): a directed self-loop $(w, w) \in E$ for some $w \in W$ would lift under item~iii.\ to the directed edge $(w^i, w^o) \in E_{\swig(W)}$, whose index values satisfy $\phi(w^i) = j + \tfrac{1}{3} > j - \tfrac{1}{3} = \phi(w^o)$, contradicting the parent-precedence requirement $\phi(\mathrm{source}) < \phi(\mathrm{target})$ of def \ref{def-topological-order}. Acyclicity of $G$ -- specifically the consequence ``no directed self-loops on $V$'' established from def \ref{def-acylic} -- rules this out: $(w, w) \notin E$ for every $w \in V$, hence for every $w \in W \ins V$.
    \end{itemize}
    Note that, in contrast to claim_3_6 for the node-splitting construction, no directed self-loop $(w, w) \in E$ for $w \in W$ can produce a length-$2$ cycle $w^o \tuh w^i \tuh w^o$ in $G_{\swig(W)}$ here, because no transfer edge $(w^o, w^i)$ is included in $E_{\swig(W)}$ by item~iii.'s typing (cf.\ def \ref{def:G_node-splitting_intervention}, closing paragraph ``Self-loops on $W$ produce no cycles in $G_{\swig(W)}$''); the self-loop is nevertheless ruled out by the CADMG hypothesis for the topological-order reason above. The bidirected-edge set $L$ plays no role in either sub-claim: acyclicity (def \ref{def-acylic}) is defined via the non-existence of non-trivial directed walks (which use only $E$), and a topological order (def \ref{def-topological-order}) is defined via the parent-set $\Pa^G$ (def \ref{def_3_5}, item~i., which is defined from $E$ alone).
\end{Rem}

\end{document}
